% 第三章 任意构型 霍普夫子 的理论模型与构建
\chapter{任意构型 霍普夫子 的理论模型与程序化构建}

\section{霍普夫子 理论模型回顾}

微磁学模拟的能量弛豫收敛程度高度依赖于初始磁化分布的精确性。对于拓扑荷 $Q_H=1$ 的基础霍普夫子构型，通常可采用 Sutcliffe 模型的解析解进行初始化。但针对高阶拓扑荷或具备复杂空间旋转角度的磁构型，该模型不再适用，需引入基于环面坐标系的广义理论框架\cite{guslienko2025magnetic}。为解决上述框架难以直接应用于数值计算的问题，本章依据局域旋转矩阵与环面坐标映射理论，开发了一种可被微磁学求解器调用的程序模块。该方案支持任意拓扑荷与反铁磁（AFM）交替背景的离散化生成，同时提供配套的三维等值面可视化与解调算法，为后续各章的数值仿真构建技术基础。

\subsection{磁化场的旋转矩阵表示}
连续场理论将 霍普夫子 定义为嵌入均匀磁化背景的三维拓扑孤子。若取背景磁化 $\mathbf{m}_0 = \hat{z}$,空间中任意构型 $\mathbf{m}(\mathbf{r})$ 均可经坐标相关的局部旋转矩阵 $\hat{R}(\mathbf{r})$ 由背景生成,即 $\mathbf{m}(\mathbf{r}) = \hat{R}(\mathbf{r}) \mathbf{m}_0$。已有工作\cite{guslienko2025magnetic}在此基础上建立了完整的体系，局域旋转被分解为绕对称轴转角与旋转轴方位角两个独立自由度,由此可严格证明 霍普夫子 拓扑荷等于环面坐标下极向缠绕数与方位向缠绕数之积,并导出以旋转角剖面函数为核心变量的交换能密度解析表达。

为建立从 $R^3$ 到 $S^3$ 的映射,引入四元数参数化向量 $n_\alpha = (n_1, n_2, n_3, n_4)$,各分量满足归一化条件 $\sum n_\alpha^2 = 1$。局部旋转矩阵据此构造为:
\begin{equation}
    R_{ik} = \delta_{ik} + 2(n_i n_k - n^2 \delta_{ik}) - 2\varepsilon_{ikl} n_l n_4
\end{equation}
其中,下标 $i, k, l \in \{1,2,3\}$,$\mathbf{n} = (n_1, n_2, n_3)$,$n^2 = n_1^2 + n_2^2 + n_3^2$,$\varepsilon_{ikl}$ 为完全反对称张量。将上式代入背景磁化 $\mathbf{m}_0 = \hat{z}$ 并展开,三维磁化场各分量可归结为:
\begin{equation}
    \mathbf{m} = (1 - 2n^2)\hat{z} + 2n_z \mathbf{n} + 2n_4(\mathbf{n} \times \hat{z})
    \label{eq:hopf_magnetization}
\end{equation}

\subsection{环面坐标下的参数化}
直角坐标 $(x,y,z)$ 下对 霍普夫子 核心形貌的解析描述牵涉大量几何约束,处理效率不高。环面坐标系 $(\eta, \beta, \phi)$ 则提供了一条更为简洁的途径，$R^3$ 空间到超球面 $S^3$ 的映射在此坐标下自然成立。两套坐标间的变换关系如下:
\begin{equation}
    \rho = \frac{a \sinh \eta}{\tau}, \quad z = \frac{a \sin \beta}{\tau}, \quad \phi = \arctan(y/x)
\end{equation}
其中 $\rho = \sqrt{x^2+y^2}$,$\tau = \cosh \eta - \cos \beta$,$a$ 为 霍普夫子 核心环($m_z=-1$ 所在的圆环)的特征大半径。

非平凡 Hopf 映射的构造要求在该坐标系下引入如下一组复函数:
\begin{align}
    Z_1 &= n_4 - i n_3 = e^{-im\beta} \cos \zeta(\eta) \\
    Z_2 &= n_2 - i n_1 = e^{in\phi} \sin \zeta(\eta)
\end{align}
将上述定义代入\cref{eq:hopf_magnetization},磁化矢量的面外分量与面内分量分别简化为:
\begin{align}
    m_z(\eta, \beta, \phi) &= \cos(2\zeta(\eta)) \\
    m_x + i m_y &= \sin(2\zeta(\eta)) e^{i(n\phi + m\beta)}
\end{align}
式中整数 $p=m$ 与 $q=n$ 分别对应极向和方位向缠绕数,系统总 Hopf 指数由两者之积唯一给定,即 $Q_H = p \times q$。衰减函数 $\zeta(\eta)$ 控制从核心到边界的径向过渡特性,其边界条件负责约束中心区和远场的磁化渐进行为。上述理论体系将高维拓扑构型的完整信息浓缩为有限个整数参数，为程序化实现提供了天然接口。本文后续工作的核心即在于将这组解析表达式转化为可在三维离散网格上执行的生成算法。

\section{离散化生成脚本与反铁磁拓展}

\subsection{程序化离散与网格构建}
以上节解析表达式为蓝本,本文采用 Python 语言开发了磁化场初始态的程序化生成工具。算法流程可拆解为三个阶段，首先在预设的三维计算域中构建均匀间距的离散网格,获取坐标矩阵 $(x_i, y_j, z_k)$；对坐标施加中心偏移与平移校正后,逐点完成从直角坐标到局部环面坐标的转换；最终在每一格点处按解析公式计算磁化三分量 $\mathbf{m}(x_i, y_j, z_k)$。

核心环所在平面可通过参数开关灵活指定，默认 $x$-$y$ 平面之外,$x$-$z$ 与 $y$-$z$ 平面均可经坐标轮换变换切换到位。该坐标转换机制使拓扑磁构型能够在三维空间内沿任意维度旋转，从而满足动力学仿真中关于初始磁化取向与自旋波波矢夹角的设定需求。离散化完成后,全部磁化数据以 OOMMF OVF 2.0 二进制格式输出,可被 Mumax3 等主流微磁学求解器直接加载为初始态。

\subsection{反铁磁背景}
针对反铁磁交换耦合体系的研究需求，初始化脚本中增设了反铁磁（AFM）背景控制模块。该模块沿各空间维度叠加符号交替矩阵,支持三维棋盘格与沿指定轴向的分层两种模式,从而可将多种反铁磁序叠加于已构建的铁磁 霍普夫子 之上。开关参数进一步提供了作用域控制，交替扰动既可仅施加于 霍普夫子 核心区域之外,亦可扩展至整个计算域以实现全局交错排列。

\section{霍普夫子 构型的可视化}

构造算法的正确性要通过可视化手段来验证。为此,本文开发了一套集等值面提取与三维渲染于一体的分析脚本,用于直观呈现 霍普夫子 的拓扑几何特征。

脚本的核心流程以行进立方体（Marching Cubes）算法从离散网格数据中抽取 $m_z=0$ 等值面,再依据面上各顶点的面内相位角 $\Phi = \arctan(m_y/m_x)$ 赋予色相值。通过色相映射(Hue Map),磁化矢量在几何曲面上的连续扭转得以直观呈现。

从仿真结果来看, AFM 交替背景的结构文件中,网格级正负磁矩的逐点翻转会导致常规等值面提取算法失效。针对这一问题,可视化脚本内置了自动信号解调模块。该模块首先利用相邻格点的平均点积识别反铁磁分布模式及其相位偏移,再等值面提取之前对局部反向磁矩执行预翻转,将场分布还原为连续的铁磁形式。解调完成后,等值面可正常渲染,核心大半径 $R$ 与管半径 $r$ 的定量提取也随之实现。

图~\ref{fig:hopfion_vis} 给出四种代表性构型的三维渲染。(a) 对应最基础的 $Q_H=1$ 构型；(b) 和 (c) 的拓扑荷同为 $Q_H=2$,却分属 $(p=2, q=1)$ 与 $(p=1, q=2)$ 两种缠绕数组合,形态呈现出明显区分；(d) 则为 $Q_H=4$ 高阶构型。色相沿等值面平滑渐变,反映出面内磁化相位的连续梯度特征。几何拓扑与矢量场分布两方面的一致性,为环面坐标理论模型以及本文离散化生成程序的可靠性提供了交叉验证。

\begin{figure}[H]
    \centering
    \includegraphics[width=0.75\textwidth]{figures/fig3-1_hopfion_vis_combined.png}
    \caption{四种拓扑构型 霍普夫子 的结构图。(a) $Q_H=1$, $p=1, q=1$; (b) $Q_H=2$, $p=2, q=1$; (c) $Q_H=2$, $p=1, q=2$; (d) $Q_H=4$, $p=2, q=2$。其中 (b) 与 (c) 具有相同的 Hopf 指数但不同的缠绕数组合。理论特征尺寸统一设定为 $R=\SI{12}{nm}, r=\SI{6}{nm}$。}
    \label{fig:hopfion_vis}
\end{figure}

本章针对任意拓扑荷霍普夫子初始态的构建问题，提出了一种基于理论解析表达式的离散化生成方法。理论计算环节将四元数体系与环面坐标映射转化为离散格点数据。程序实现环节输出的 OVF 格式场分布覆盖了不同 $(p,q)$ 缠绕数参数空间，并集成了自适应反铁磁背景生成功能。等值面可视化脚本借助自动 AFM 解调与面内相位着色机制，分别从几何形貌和矢量分布两个角度对生成结果实施了校核。相较于针对单一构型的解析方案，该离散化生成工具扩展了初始态拓扑荷的取值范围，为后续的阻挫体系稳定性分析与自旋波动力学模拟提供了前置条件。
